% Draft sections: Introduction + Related Work.
% This file is content-only and is intended to be \input{} into docs/cmc/main.tex.

\section{Introduction}
\label{sec:introduction}

Automatic surface inspection is a persistent requirement in industrial production because manual inspection is costly, subjective, and difficult to scale. In wood processing, the problem is even harder than in many manufactured materials because the substrate is inherently heterogeneous: grain, resin, knots, cracks, pith, stains, and growth patterns create large appearance variation even before true defects are considered. This challenge belongs to a broader wood-vision research landscape that spans species recognition, texture analysis, grading, and material characterization, where strong appearance variability, limited labeled data, and domain-specific acquisition settings have long been recognized as central obstacles \citep{hwang_sugiyama_review,da_silva_2017,souza_2020,de_geus_2021,kirbas_2022,he_2021_woodresearch,herrera_poyatos_2024}. Kodytek \textit{et al.} note that surface defects do not merely lower visual quality; they can also degrade strength and stiffness, which directly affects material utilization and economic value in downstream processing \citep{kodytek_dataset}. In practice, an inspection system must therefore detect visually weak defects, separate semantically similar subtypes, and remain stable across production batches while using detector choices that are realistic for deployment rather than only optimal on a curated benchmark.

The defect types themselves are also operationally meaningful. The large-scale dataset considered in our study contains ten common categories, including live knots, dead knots, quartzity, knots with cracks, missing knots, cracks, overgrown regions, resin, marrow, and blue stain \citep{kodytek_dataset}. A live knot usually corresponds to a branch remnant that remains relatively intergrown with surrounding fibers, whereas a dead knot is more detached and may indicate weaker local structure or later breakout during machining. A knot with crack combines fiber distortion with fracture-like splitting, and a missing knot leaves a cavity or void that directly reduces usable surface quality. Longitudinal cracks can propagate during drying, handling, or cutting. Resin creates contamination and finishing problems, marrow often lowers appearance grade and dimensional stability, and blue stain reduces commercial value even when the mechanical damage is less obvious. In the present study, we focus on seven frequent and practically important defect classes: live knot, dead knot, knot with crack, knot missing, crack, resin, and marrow. Representative images from the source benchmark and the low-resource target domain are compared later in Section~\ref{sec:datasets}.

From a vision perspective, wood surface defect detection is challenging for at least four reasons. First, many targets are small relative to the board image and may appear low-contrast, elongated, or only partially visible, which makes localization unstable under repeated downsampling. Second, several categories such as live knots, dead knots, knots with cracks, and cracks are semantically close and visually entangled, so subtype confusion is easy even when a detector fires near the right region. Third, the class distribution is long-tailed, which makes rare but operationally meaningful defects hard to learn. Fourth, deployment data rarely match benchmark data in framing, scale, label taxonomy, or defect prevalence, so transfer behavior matters as much as in-domain validation.

The research response to these problems has been dominated by detector redesign. Representative wood-defect models include ODCA-YOLO \citep{odca_yolo}, BPN-YOLO \citep{bpn_yolo}, SiM-YOLO \citep{sim_yolo}, FDD-YOLO \citep{fdd_yolo}, and CWB-YOLOv8 \citep{an_cwb_yolov8}. Related wood-inspection studies have also explored anomaly-oriented detection and hybrid sensing setups \citep{kilic_2024,kilic_2025_wd_detector}. Comparable design priorities also appear in recent inspection studies on faulty wood regions, difficult small objects, PCB faults, rail defects, and metallic or steel surfaces \citep{cmc_wood_area_attention,cmc_insulator_small,cmc_pcb_yolov8,cmc_cg_yolov8,cmc_rail_upernet,cmc_yolo_rlc,cmc_lmvit,cmc_msese,cmc_flame,cmc_channel_shuffle}. Taken together, these works show that small-object sensitivity, multiscale representation, localization stability, and deployment efficiency remain recurring pain points across defect domains.

Despite these advances, three practical questions remain insufficiently resolved. First, it is still unclear whether improvements reported on screened benchmarks reflect genuine detector superiority or protocol-specific optimization. Second, direct comparisons across detector families under matched conditions remain limited, even though such comparisons are essential for deployment-oriented model selection. Third, the extent to which in-domain gains transfer to external datasets remains underexplored. These unresolved questions matter because wood inspection systems are expected to operate under constrained resources, heterogeneous imaging conditions, and shifting target distributions.

Building upon this problem setting, our work addresses a question that remains underexamined in the current literature: how detector family, evaluation protocol, and dataset shift jointly shape observed performance in wood surface defect detection. We organize the study in three stages. First, we examine lightweight Faster R-CNN baselines using MobileNetV3-FPN backbones, including a low-capacity exploratory baseline and a higher-resolution \texttt{MobileNet\_HR} variant. Second, we compare these results with a compact YOLOv8 baseline centered on YOLOv8s. Third, after identifying YOLOv8s as the strongest compact in-domain model, we test two plausible refinements: a P2 small-object head and a WNIoU-style hybrid box loss.

This framing is motivated by a practical observation: a detector that appears strong only on a screened benchmark, or a refinement that improves only one local setting, offers limited deployment value. Our study therefore treats performance differences across protocols and datasets not as secondary details, but as central evidence for judging model usefulness in practice.

This study makes the following contributions.

\begin{enumerate}[leftmargin=1.5em]
  \item We present a comparative study of wood surface defect detection that connects results on a public screened benchmark with performance under a broader in-domain protocol and a low-resource target-domain setting.
  \item We compare lightweight two-stage detectors and a compact one-stage detector family under shared data processing and a common internal evaluation pipeline where possible, showing that detector-family choice has a larger effect than the tested local architectural refinements in the current task.
  \item We analyze benchmark sensitivity by combining study-internal 20-epoch screening comparisons with a 200-epoch literature-facing YOLO confirmation run, showing that the screened benchmark remains more favorable than the broader in-domain protocol and that an additional 200-epoch full-protocol rerun does not materially close this gap under the current setup.
  \item We analyze target-domain adaptation on a low-resource wood-inspection dataset and show that source-domain pretraining improves supervised adaptation under a fixed 50-epoch budget, while domain shift remains a substantial challenge.
\end{enumerate}

The remainder of this paper is organized as follows. Section~\ref{sec:related-work} reviews recent wood-defect and industrial-defect detection literature, with emphasis on lightweight detectors, YOLO-style baselines, and small-object-oriented refinements. Section~\ref{sec:methods} describes the datasets, protocol variants, detector families, and evaluation setup. Section~\ref{sec:experiments} presents the in-domain and external experiments. Section~\ref{sec:discussion} discusses the main findings, limitations, and practical implications.

\section{Related Work}
\label{sec:related-work}

\subsection{Industrial Surface Inspection and Practical Detector Constraints}

Recent research on automated wood and industrial defect inspection can be organized into four broad lines: lightweight two-stage detectors, compact YOLO-style one-stage detectors, small-object-oriented detector refinements, and benchmarking studies that emphasize protocol sensitivity or generalization. This grouping is consistent with recent reviews of industrial defect inspection, edge-oriented detection, and small-object analysis, and it also matches the design space explored in our study \citep{ma_2024_air,wang_2025_ieeeaccess,mittal_2024,wei_2024_smallreview,hua_2025_smallsurvey}.

What these reviews make clear is that industrial inspection is no longer judged by accuracy alone. In practice, a detector must satisfy latency and hardware constraints, cope with class imbalance, and remain usable when the operating data differ from the development benchmark. This is one reason why two-stage detectors such as Faster R-CNN remain attractive as controlled references: they support explicit proposal refinement and flexible backbone substitution \citep{faster_rcnn}. At the same time, lightweight backbones such as MobileNetV3 are widely favored when resource usage matters, while YOLO-style one-stage detectors have become the default practical choice because they offer dense prediction with favorable speed--accuracy trade-offs \citep{mobilenetv3,mittal_2024}.

Recent studies on printed-surface inspection, metallic and steel surface defects, and rail defect analysis likewise emphasize efficiency, robustness, or multiscale enhancement rather than only raw benchmark gains \citep{cmc_printed_positive,cmc_channel_shuffle,cmc_lmvit,cmc_msese,cmc_rail_upernet}. The recurring limitation is not a shortage of detector ideas, but a shortage of detector-selection evidence under realistic protocol shifts. Many papers show that a modified detector improves one benchmark, but fewer ask whether a lightweight two-stage family, a compact one-stage family, and a scaled-up model behave differently once the protocol broadens or the data source changes.

\subsection{Wood Surface Defect Detection on Screened Benchmarks}

The most directly relevant literature concerns wood surface defect detection with one-stage detectors on curated subsets of public or self-constructed datasets. The large-scale VSB dataset of Kodytek \textit{et al.} has become particularly influential because it offers a clearer defect taxonomy and a more systematic benchmark than many earlier in-house collections \citep{kodytek_dataset}. Recent detector papers built on this line of work follow a familiar pattern: they keep a YOLO-style backbone but introduce custom fusion, attention, small-target handling, or localization losses to raise the headline score. This is the basic trajectory seen in ODCA-YOLO, BPN-YOLO, SiM-YOLO, FDD-YOLO, and the more recent CWB-YOLOv8 study \citep{odca_yolo,bpn_yolo,sim_yolo,fdd_yolo,an_cwb_yolov8}.

There is also a neighboring stream of wood-inspection work that is less tied to standard bounding-box detection. Some studies approach the problem through anomaly detection, hybrid sensing, or region-oriented quality analysis rather than direct multi-class detection \citep{kilic_2024,kilic_2025_wd_detector,cmc_wood_area_attention}. Taken together, these studies suggest three points. First, YOLO-like one-stage detectors have become the dominant practical backbone for wood defect detection. Second, most reported gains still come from module-level architectural tailoring on screened or curated settings, with headline scores clustered in a relatively narrow range around the low-80s. Third, comparisons across detector families, protocol breadth, and external transfer remain much rarer than improvements reported within a single benchmark setup.

\subsection{Lightweight and Resource-Aware Defect Detectors}

The practical motivation for lightweight models is especially visible in recent defect-detection work. Reviews of edge-oriented object detection repeatedly emphasize backbone efficiency, memory footprint, and inference cost as decisive factors in deployment, especially when models are expected to run close to the acquisition pipeline rather than on unconstrained offline servers \citep{mittal_2024}. The same deployment pressure is visible in compact PCB, aluminum, printed-surface, and metallic-surface detectors, where authors prioritize lightweight operators, efficient heads, or reduced-complexity backbones while still trying to preserve acceptable inspection accuracy \citep{cmc_pcb_yolov8,cmc_yolo_rlc,cmc_cg_yolov8,cmc_aluminum_yolov8n,cmc_printed_positive,cmc_channel_shuffle}. 

What is still missing from this literature is a clearer answer to a basic engineering question: when should one keep optimizing a lightweight family, and when is it better to move to a stronger compact detector? Most studies stay inside one detector family, so they reveal how to improve that family but not whether another family would be a better practical choice under the same protocol.

This gap matters in wood inspection. If a MobileNet-based two-stage detector is already close to a compact YOLO baseline, then the choice may favor a different engineering trade-off than if the gap is large. Conversely, if the stronger detector only wins on one screened subset, then the gain may not be operationally meaningful. Our work addresses this specific comparison directly.

\subsection{Small-Object and Defect-Oriented Detector Modifications}

Beyond wood inspection, industrial-defect research repeatedly returns to the same three bottlenecks: weak small-object representation, strong multiscale variation, and unstable localization. Small-object surveys consistently point to high-resolution heads, feature-pyramid redesign, attention-guided fusion, and alternative localization losses as the most common responses to these issues \citep{wei_2024_smallreview,hua_2025_smallsurvey}. The same pattern appears across recent studies of insulator, rail, PCB, aluminum, and steel defects, where higher-resolution prediction layers, multiscale fusion, enhanced receptive fields, or revised loss formulations are used to stabilize localization under difficult imaging conditions \citep{cmc_insulator_small,cmc_yolovsi,cmc_cg_yolov8,cmc_pcb_yolov8,cmc_aluminum_yolov8n,cmc_msese}. More representation-oriented approaches, such as learnable-memory ViTs and federated unsupervised pipelines, broaden the design space further \citep{cmc_lmvit,cmc_flame}.

This literature is important for our study because it shows that the tested refinements, especially the P2 head and the modified regression loss, are well grounded in current practice rather than being ad hoc additions. At the same time, it exposes a persistent weakness: architecture-first improvements are usually validated on a single dataset and a single protocol, which makes it difficult to distinguish genuinely robust gains from gains that mainly fit one curated setup \citep{bpn_yolo,sim_yolo,fdd_yolo,an_cwb_yolov8}.

\subsection{Benchmarking, Reproducibility, and Transfer Gaps}

The final related-work thread concerns evaluation itself. Recent benchmarking work on surface defect detection argues that reproducibility, statistical rigor, and protocol transparency deserve much more attention \citep{defect_benchmark_rigorous}. Broader reviews of industrial defect inspection make the same point more generally, noting that inconsistent datasets, metrics, and acquisition conditions still make cross-paper comparison difficult \citep{ma_2024_air,wang_2025_ieeeaccess}. This concern is especially relevant to wood surface inspection, where benchmark curation, class filtering, and screening choices can materially change the apparent ranking of detectors. A detector may look strong on a screened subset yet degrade under a broader protocol, a different sampling scheme, or an external dataset with a partially mismatched label space. From a deployment standpoint, such failure modes are not a side issue; they largely define whether the system is usable.

Building upon this research context, our work differs from the dominant architecture-centric pattern in one key respect: we do not introduce another detector block. Instead, we organize the study around four practical questions: how lightweight MobileNet-based two-stage baselines compare with compact YOLO baselines, how stable the compact YOLO baseline remains when moving from a screened benchmark to a broader full seven-class protocol, whether common small-object-oriented refinements actually help that compact baseline during the screening stage, and whether a screened source model remains useful when adapting to a target-domain dataset under a limited training budget. In our view, the main weakness of much prior work is not that the proposed modules are unreasonable; it is that they are usually evaluated under conditions too narrow to answer these broader deployment-facing questions.

\subsection{Position of the Present Work}

In summary, prior work suggests three relevant solution families for the present study: lightweight backbone selection, compact one-stage detectors, and small-object-oriented detector refinements. Our experiments cover all three directions in a structured comparative setting. Faster R-CNN + MobileNetV3 baselines represent the lightweight family, YOLOv8s represents the compact one-stage family, and \texttt{Y1} and \texttt{Y2} represent head- and loss-level refinements. We use these families to identify which conclusions remain stable once benchmark sensitivity and target-domain adaptation are incorporated into the evaluation. This is the specific gap addressed by the present study.
